\FloatBarrier
\section{Esercizio Guidato: Matching Completo ed Enumerazione delle Istanze}
\label{sec:dm-seq-exercise}

Per padroneggiare la logica dell'algoritmo di matching, svolgiamo un esercizio completo di verifica e censimento di tutte le istanze valide.

\FloatBarrier
\subsection{Testo dell'Esercizio}
Sia data la sequenza di dati complessa $s$ composta da 8 elementi distribuiti da $t = 0$ a $t = 7$:
\begin{equation}
\label{eq:seq-exercise-input}
\begin{split}
s = \langle &\{A, C\}_{t=0} \, \{C, D\}_{t=1} \, \{F, H\}_{t=2} \, \{A, B\}_{t=3} \\
&\{B, C, D\}_{t=4} \, \{E\}_{t=5} \, \{A, B, D\}_{t=6} \, \{F\}_{t=7} \rangle
\end{split}
\end{equation}

Si determini se i seguenti quattro pattern candidati sono contenuti in $s$, e in caso affermativo si enumerino tutte le istanze distinte $\langle t_1, \dots, t_n \rangle$:
\begin{enumerate}
    \item $P_1 = \langle \{C\} \, \{H\} \, \{C\} \rangle$
    \item $P_2 = \langle \{A\} \, \{F\} \rangle$
    \item $P_3 = \langle \{A\} \, \{A\} \, \{D\} \rangle$
    \item $P_4 = \langle \{A\} \, \{A, B\} \, \{F\} \rangle$
\end{enumerate}

\FloatBarrier
\subsection{Risoluzione Analitica Dettagliata}

\subsubsection{Pattern 1: $P_1 = \langle \{C\} \, \{H\} \, \{C\} \rangle$}
Il pattern richiede di individuare tre marcatori temporali $t_1 < t_2 < t_3$ tali che:
\begin{equation*}
C \in s_{t_1}, \quad H \in s_{t_2}, \quad C \in s_{t_3}
\end{equation*}
\begin{itemize}
    \item L'elemento $\{H\}$ è presente unicamente al tempo $t = 2$ ($s_2 = \{F, H\}$). Pertanto deve valere necessariamente $t_2 = 2$.
    \item Cerchiamo le occorrenze di $C$ per $t_1 < 2$:
    \begin{itemize}[noitemsep]
        \item $t_1 = 0$: $C \in s_0 = \{A, C\}$;
        \item $t_1 = 1$: $C \in s_1 = \{C, D\}$.
    \end{itemize}
    \item Cerchiamo le occorrenze di $C$ per $t_3 > 2$:
    \begin{itemize}[noitemsep]
        \item Al tempo $t = 4$: $C \in s_4 = \{B, C, D\}$. Nei tempi successivi ($5, 6, 7$) l'item $C$ non compare più. Dunque $t_3 = 4$.
    \end{itemize}
\end{itemize}
Combinando le scelte possibili, otteniamo esattamente \textbf{2 istanze distinte}:
\begin{enumerate}
    \item $\langle t_1 = 0, \, t_2 = 2, \, t_3 = 4 \rangle$
    \item $\langle t_1 = 1, \, t_2 = 2, \, t_3 = 4 \rangle$
\end{enumerate}
\textbf{Esito}: Il pattern $P_1$ \textbf{OCCORRE} (è supportato con 2 istanze).

\subsubsection{Pattern 2: $P_2 = \langle \{A\} \, \{F\} \rangle$}
Il pattern richiede di individuare due istanti $t_1 < t_2$ tali che:
\begin{equation*}
A \in s_{t_1} \quad \text{e} \quad F \in s_{t_2}
\end{equation*}
\begin{itemize}
    \item L'item $A$ è presente nei tempi: $t \in \{0, 3, 6\}$ (poiché $s_0 = \{A, C\}$, $s_3 = \{A, B\}$, $s_6 = \{A, B, D\}$).
    \item L'item $F$ è presente nei tempi: $t \in \{2, 7\}$ (poiché $s_2 = \{F, H\}$, $s_7 = \{F\}$).
\end{itemize}
Valutiamo tutte le coppie $(t_1, t_2)$ con $t_1 < t_2$:
\begin{itemize}
    \item Per $t_1 = 0$: i valori ammissibili per $t_2 > 0$ sono $t_2 = 2$ e $t_2 = 7 \implies \langle 0, 2 \rangle$ e $\langle 0, 7 \rangle$.
    \item Per $t_1 = 3$: l'unico valore per $t_2 > 3$ è $t_2 = 7 \implies \langle 3, 7 \rangle$.
    \item Per $t_1 = 6$: l'unico valore per $t_2 > 6$ è $t_2 = 7 \implies \langle 6, 7 \rangle$.
\end{itemize}
Otteniamo esattamente \textbf{4 istanze distinte}:
\begin{equation}
\langle 0, 2 \rangle, \quad \langle 0, 7 \rangle, \quad \langle 3, 7 \rangle, \quad \langle 6, 7 \rangle
\end{equation}
\textbf{Esito}: Il pattern $P_2$ \textbf{OCCORRE} (è supportato con 4 istanze).

\subsubsection{Pattern 3: $P_3 = \langle \{A\} \, \{A\} \, \{D\} \rangle$}
Il pattern richiede tre istanti $t_1 < t_2 < t_3$ con:
\begin{equation*}
A \in s_{t_1}, \quad A \in s_{t_2}, \quad D \in s_{t_3}
\end{equation*}
\begin{itemize}
    \item Le posizioni contenenti $A$ sono $t \in \{0, 3, 6\}$. Le coppie ordinate $(t_1, t_2)$ con $t_1 < t_2$ sono:
    \begin{equation*}
    (0, 3), \quad (0, 6), \quad (3, 6)
    \end{equation*}
    \item L'item $D$ compare nei tempi: $t \in \{1, 4, 6\}$ (poiché $s_1 = \{C, D\}$, $s_4 = \{B, C, D\}$, $s_6 = \{A, B, D\}$).
    \item Verifichiamo ciascuna coppia $(t_1, t_2)$:
    \begin{itemize}
        \item Per $(t_1 = 0, t_2 = 3)$: cerchiamo $t_3 > 3$ con $D \in s_{t_3}$. I tempi validi sono $t_3 = 4$ e $t_3 = 6$.
        Ciò genera le istanze $\langle 0, 3, 4 \rangle$ e $\langle 0, 3, 6 \rangle$.
        \item Per $(t_1 = 0, t_2 = 6)$: richiederebbe $t_3 > 6$ con $D \in s_{t_3}$. Al tempo $t = 7$ è presente solo $\{F\}$, quindi nessun $D$ è disponibile dopo $t = 6$. Nessuna istanza valida.
        \item Per $(t_1 = 3, t_2 = 6)$: analogamente richiederebbe $t_3 > 6$, impossibile per assenza di $D$ a $t > 6$. Nessuna istanza valida.
    \end{itemize}
\end{itemize}
Otteniamo esattamente \textbf{2 istanze distinte}:
\begin{equation}
\langle 0, 3, 4 \rangle \quad \text{e} \quad \langle 0, 3, 6 \rangle
\end{equation}
\textbf{Esito}: Il pattern $P_3$ \textbf{OCCORRE} (è supportato con 2 istanze).

\subsubsection{Pattern 4: $P_4 = \langle \{A\} \, \{A, B\} \, \{F\} \rangle$}
Il pattern richiede tre istanti $t_1 < t_2 < t_3$ tali che:
\begin{equation*}
A \in s_{t_1}, \quad \{A, B\} \subseteq s_{t_2}, \quad F \in s_{t_3}
\end{equation*}
\begin{itemize}
    \item L'itemset congiunto $\{A, B\}$ è contenuto in:
    \begin{itemize}[noitemsep]
        \item $s_3 = \{A, B\} \implies t_2 = 3$;
        \item $s_6 = \{A, B, D\} \implies t_2 = 6$.
    \end{itemize}
    \item \textbf{Caso $t_2 = 3$}:
    \begin{itemize}[noitemsep]
        \item Per $t_1 < 3$ con $A \in s_{t_1}$: l'unica scelta è $t_1 = 0$ ($s_0 = \{A, C\}$).
        \item Per $t_3 > 3$ con $F \in s_{t_3}$: l'unica scelta è $t_3 = 7$ ($s_7 = \{F\}$).
        \item Istanza generata: $\langle 0, 3, 7 \rangle$.
    \end{itemize}
    \item \textbf{Caso $t_2 = 6$}:
    \begin{itemize}[noitemsep]
        \item Per $t_1 < 6$ con $A \in s_{t_1}$: possiamo scegliere $t_1 = 0$ oppure $t_1 = 3$.
        \item Per $t_3 > 6$ con $F \in s_{t_3}$: l'unica scelta è $t_3 = 7$.
        \item Istanze generate: $\langle 0, 6, 7 \rangle$ e $\langle 3, 6, 7 \rangle$.
    \end{itemize}
\end{itemize}
Otteniamo esattamente \textbf{3 istanze distinte}:
\begin{equation}
\langle 0, 3, 7 \rangle, \quad \langle 0, 6, 7 \rangle, \quad \langle 3, 6, 7 \rangle
\end{equation}
\textbf{Esito}: Il pattern $P_4$ \textbf{OCCORRE} (è supportato con 3 istanze).

\FloatBarrier
\subsection{Quadro Riepilogativo delle Risoluzioni}
La Tabella~\ref{tab:dm-seq-exercise-summary} sintetizza i risultati dell'analisi per la sequenza $s$.

\begin{table}[H]
\centering
\small
\renewcommand{\arraystretch}{1.25}
\begin{tabularx}{\textwidth}{c Y c Y}
\toprule
\textbf{Pattern} & \textbf{Forma Strutturale} & \textbf{Esito} & \textbf{Istanze Temporali $\langle t_1, \dots, t_n \rangle$} \\
\midrule
$P_1$ & $\langle \{C\} \, \{H\} \, \{C\} \rangle$ & \textbf{OCCORRE} & $\langle 0, 2, 4 \rangle, \; \langle 1, 2, 4 \rangle$ (2 istanze) \\
$P_2$ & $\langle \{A\} \, \{F\} \rangle$ & \textbf{OCCORRE} & $\langle 0, 2 \rangle, \; \langle 0, 7 \rangle, \; \langle 3, 7 \rangle, \; \langle 6, 7 \rangle$ (4 istanze) \\
$P_3$ & $\langle \{A\} \, \{A\} \, \{D\} \rangle$ & \textbf{OCCORRE} & $\langle 0, 3, 4 \rangle, \; \langle 0, 3, 6 \rangle$ (2 istanze) \\
$P_4$ & $\langle \{A\} \, \{A, B\} \, \{F\} \rangle$ & \textbf{OCCORRE} & $\langle 0, 3, 7 \rangle, \; \langle 0, 6, 7 \rangle, \; \langle 3, 6, 7 \rangle$ (3 istanze) \\
\bottomrule
\end{tabularx}
\caption{Quadro riassuntivo delle istanze e del matching per i quattro pattern dell'Esercizio.}
\label{tab:dm-seq-exercise-summary}
\end{table}

\begin{noteblock}{Sintesi Metodologica}
    L'analisi condotta conferma come:
    \begin{enumerate}[noitemsep]
        \item Un singolo elemento della sequenza candidato possa contenere uno o più item simultanei;
        \item La sequenzialità temporale richieda strettamente indici crescenti ($t_1 < t_2 < \dots < t_n$);
        \item Elementi con item identici possano ripetersi in istanti temporali distinti all'interno della stessa sequenza;
        \item Nel modello standard di Sequential Pattern Mining, ciascuno di questi pattern incrementerebbe di $+1$ il conteggio di supporto di questa sequenza, indipendentemente dalla molteplicità degli embedding rilevati.
    \end{enumerate}
\end{noteblock}
